% =====================================================================
%  統計基礎 1. 確率と確率分布
%  構成は Docs/SectionPlan.md にしたがう。
% =====================================================================
\chapter{確率と確率分布}
\label{chap:stat-probability}

\begin{chaptergoal}
  % TODO: この章の目標
\end{chaptergoal}

% ---------------------------------------------------------------------
\section{確率分布、確率変数}
\label{sec:stat-moments}

\keywords{チェビシェフの不等式、積率、尖度、歪度、積率母関数}

\subsection{積率（モーメント）}
% TODO:   原点のまわりの $k$ 次積率 $E[X^k]$ と、平均のまわりの $k$ 次積率 $E[(X-\mu)^k]$
% TODO:   期待値と分散を積率で表す

\subsection{歪度（わいど）}
% TODO:   3次の積率を標準化したもの
% TODO:   分布の左右のかたよりを表す。正・負・0 の場合を図で比べる

\subsection{尖度（せんど）}
% TODO:   4次の積率を標準化したもの
% TODO:   正規分布では 3 になる。3 を引いた「超過尖度」の流儀もあることに注意
% TODO:   分布のとがり具合と、裾（すそ）の重さを表す

\subsection{積率母関数}
% TODO:   定義 $M_X(t)=E[e^{tX}]$
% TODO:   微分して $t=0$ を代入すると積率が得られる
% TODO:   積率母関数が同じなら分布も同じ（分布を見分ける道具になる）
% TODO:   独立な確率変数の和の積率母関数は、それぞれの積になる
% TODO:   例：二項分布、ポアソン分布、正規分布の積率母関数
% TODO:   〔数学メモ〕$e^x$ の微分、テイラー展開のイメージ

\subsection{チェビシェフの不等式}
% TODO:   主張：$P(|X-\mu|\ge k\sigma)\le 1/k^2$
% TODO:   証明（マルコフの不等式を経由する）
% TODO:   分布の形を知らなくても使えるという利点
% TODO:   正規分布での実際の確率と比べると、評価がゆるいこと

\subsection{確率変数の変換}
% 〔補足〕
% TODO:   $Y=g(X)$ の確率密度関数の求め方
% TODO:   対数正規分布（1.2）や逆変換法（4.2）の準備

% ---------------------------------------------------------------------
\section{主要な確率分布}
\label{sec:stat-distributions}

\keywords{対数正規分布、ガンマ分布、ベータ分布、超幾何分布、負の二項分布、多変量正規分布}

\subsection{超幾何分布}
% TODO:   「戻さずに取り出す」（非復元抽出）ときの当たりの個数の分布
% TODO:   期待値と分散
% TODO:   二項分布との関係（母集団が大きいと二項分布に近づく）、有限母集団修正

\subsection{負の二項分布}
% TODO:   「$r$ 回成功するまでに失敗した回数」の分布
% TODO:   幾何分布を一般にしたもの
% TODO:   期待値と分散
% TODO:   使われる場面：ポアソン分布よりばらつきが大きいカウントデータ（過分散）

\subsection{ガンマ分布}
% TODO:   〔数学メモ〕ガンマ関数 $\Gamma(\alpha)$：広義積分による定義、$\Gamma(n)=(n-1)!$、$\Gamma(1/2)=\sqrt{\pi}$（数学基礎「1変数関数の積分法」を参照）
% TODO:   確率密度関数、形状パラメータと尺度パラメータ
% TODO:   期待値と分散
% TODO:   指数分布の和はガンマ分布になる
% TODO:   カイ二乗分布はガンマ分布の特別な場合（2.1 への準備）

\subsection{ベータ分布}
% TODO:   〔数学メモ〕ベータ関数 $B(a,b)$ と、ガンマ関数との関係
% TODO:   0 から 1 の値をとる分布で、「確率そのもの」の分布として使える
% TODO:   一様分布は特別な場合
% TODO:   期待値と、パラメータによる形の変化
% TODO:   ベイズ統計の事前分布として使うことの予告（3.1）

\subsection{対数正規分布}
% TODO:   $\log X$ が正規分布にしたがうときの $X$ の分布
% TODO:   右に長く裾をひく形。所得や株価などの例
% TODO:   期待値 $e^{\mu+\sigma^2/2}$、中央値 $e^{\mu}$

\subsection{多変量正規分布}
% TODO:   〔数学メモ〕ベクトルと行列による表し方、分散共分散行列（数学基礎「線形代数」を参照）
% TODO:   2変量正規分布の確率密度関数と相関係数
% TODO:   等高線が楕円になること
% TODO:   周辺分布・条件付き分布も正規分布になる
% TODO:   正規分布では「無相関 ⇔ 独立」が成り立つ（一般には成り立たない）
% TODO:   線形変換しても多変量正規分布になる

\subsection{まとめ}
% ：分布どうしの関係図（二項・ポアソン・正規・ガンマ・カイ二乗・ベータ など）

% ---------------------------------------------------------------------
\section{確率変数の漸近的性質}
\label{sec:stat-asymptotics}

\keywords{大数の法則、中心極限定理、確率収束、分布収束}

\subsection{確率収束}
% TODO:   定義と直感的な意味（「ずれが大きくなる確率が 0 に近づく」）

\subsection{大数の法則}
% TODO:   主張：標本平均は母平均に確率収束する
% TODO:   チェビシェフの不等式を使った証明
% TODO:   コイン投げのシミュレーションで確かめる

\subsection{分布収束}
% TODO:   定義（累積分布関数が近づく）
% TODO:   確率収束との違い（確率収束なら分布収束、逆は成り立たない）

\subsection{中心極限定理}
% TODO:   主張：標本平均を標準化すると、標準正規分布に分布収束する
% TODO:   もとの分布が正規分布でなくてよいこと
% TODO:   二項分布の正規近似
% TODO:   積率母関数を使った証明の流れ（概略）
% TODO:   注意：分散が存在しない分布（コーシー分布）では成り立たない

\begin{chaptersummary}
  % TODO: この章のまとめ
\end{chaptersummary}
